\chapter{NKAudio : faire du bruit}

Un bouton qui ne fait aucun bruit quand on clique dessus paraît cassé. C'est
irrationnel et c'est ainsi. Ce chapitre vous donne les moyens de corriger cela
en une dizaine de lignes, puis explique tout ce qu'il y a derrière — parce que
le module va beaucoup plus loin qu'un simple « jouer un fichier ».

\texttt{NKAudio} est, des cinq modules d'intégration, celui dont la façade est
la plus aboutie. Un seul en-tête suffit :

\begin{nkcode}{L'unique inclusion nécessaire}
#include "NKAudio/NkAudio.h"
\end{nkcode}

\noindent et tout vit dans le \emph{namespace} \texttt{nkentseu::audio}.

\section{Le module, et sa dépendance surprenante}

\begin{nkcode}{Kernel/Runtime/NKAudio/NKAudio.jenga:30-35}
    nkentseudependson(
        # NKMedia : decodeur Opus from-scratch (codec .opus via NkOpusCodec).
        ["NKPlatform", "NKCore", "NKMemory", "NKContainers", "NKLogger", "NKThreading", "NKFileSystem", "NKStream", "NKMedia"],
        selfexport="NKAudio",
        extra_includes=["src"],
    )
\end{nkcode}

\texttt{NKAudio} dépend de \texttt{NKMedia}, qui dépend lui-même de
\texttt{NKImage}. La raison est mince — le décodeur Opus vit dans NKMedia — mais
elle est réelle : dans l'ordre de construction, NKImage vient d'abord, puis
NKMedia, puis NKAudio. Nous enseignons l'audio avant la vidéo parce que la
façade audio est plus simple, mais retenez la direction de la flèche : c'est
l'audio qui appelle la vidéo, jamais l'inverse.

\begin{nkcode}{Kernel/Runtime/NKAudio/ — arborescence (abrégée)}
NKAudio/
  src/NKAudio/
    NkAudio.h            (1427 l.)  <- EN-TETE PUBLIC UNIQUE : tout est la
    NkAudioEngineCore.cpp           <- AudioEngine (PIMPL)
    NkAudioLoader.cpp               <- AudioLoader (WAV natif + dispatch codecs)
    NkAudioMixer.cpp · NkAudioGenerator.cpp · NkAudioAnalyzer.cpp
    NkAudioEffects.{h,cpp}          <- effets DSP concrets
    NkAudioBackends.{h,cpp}         <- WASAPI / CoreAudio / ALSA / AAudio / WebAudio / Null
    NkAudioBus.{h,cpp}              <- bus hierarchiques Master -> SFX/Music/Voice/UI
    NkAudioCapture.{h,cpp}          <- MICRO (entree)
    Codecs/  MP3/ · OGG/ · FLAC/ · Opus/
    Streaming/
      NkAudioStream.{h,cpp}         <- IAudioStream (pull) + WavStream/MemoryStream
      NkAudioStreamPlayer.{h,cpp}   <- thread decodeur + ring buffer
      NkContainerAudioStream.{h,cpp}<- piste audio d'un conteneur video
\end{nkcode}

\section{L'invariant fondateur : un seul format interne}

Avant toute chose, comprenez ce que le module fait de vos fichiers :

\begin{nkcode}{Kernel/Runtime/NKAudio/src/NKAudio/NkAudio.h:258-260}
 @note Format interne normalisé : Float32 interleaved stéréo.
       La conversion est faite au chargement.
\end{nkcode}

Peu importe que votre fichier soit un WAV 8 bits mono à 22\,050~Hz ou un FLAC
24 bits à 96\,000~Hz : après \texttt{AudioLoader::Load}, vous avez des flottants
32 bits entrelacés. Cela simplifie énormément la suite — le mixeur n'a qu'un seul
cas à traiter — et cela signifie qu'un fichier chargé occupe souvent plus de
mémoire que sur le disque.

\begin{nkcode}{Kernel/Runtime/NKAudio/src/NKAudio/NkAudio.h:262-284}
    struct AudioSample {
            float32 *data = nullptr;   ///< Samples Float32 (interleaved)
            usize frameCount = 0;      ///< Nombre de frames (samples/channels)
            int32 sampleRate = 48000;
            int32 channels = 2;
            AudioFormat format = AudioFormat::UNKNOWN;
            float32 GetDuration() const;
            usize   GetSampleCount() const;
            bool    IsValid() const;
            memory::NkAllocator *mAllocator = nullptr; ///< Allocateur responsable
    };
\end{nkcode}

Notez le vocabulaire, qu'il faut fixer une fois pour toutes : une
\emph{\texttt{frame}} est un instant, un \emph{\texttt{sample}} est une valeur.
En stéréo, une \emph{frame} contient deux \emph{samples}. \texttt{frameCount}
compte les instants, pas les nombres — d'où
\texttt{GetSampleCount() = frameCount * channels}.

Notez aussi \texttt{mAllocator} : chaque \texttt{AudioSample} se souvient de qui
l'a alloué. Nous verrons pourquoi c'est vital.

\section{Le patron canonique en cinq temps}

L'application \texttt{NkAudioPlayer} ouvre son fichier principal sur la recette,
et il n'y a rien à y ajouter :

\begin{nkcode}{Applications/NkAudioPlayer/src/main.cpp:2-9}
// Le patron CORRECT pour lire un fichier audio vers les haut-parleurs, ET l'afficher :
//   1. AudioEngine::Instance().Initialize()  -> ouvre le device + thread de mixage
//   2. AudioLoader::Load(path)               -> décode WAV/MP3/OGG/FLAC/Opus en AudioSample
//   3. engine.Play(sample, params)           -> lance une voix (bus "Music")
//   4. boucle : GetPlaybackPosition -> tête de lecture sur la FORME D'ONDE dessinée
//   5. engine.Shutdown()
\end{nkcode}

\subsection{Temps 1 — initialiser le moteur}

\begin{nkcode}{Applications/NkAudioPlayer/src/main.cpp:113-121}
    // ── 1) Moteur audio ───────────────────────────────────────────────────────
    audio::AudioEngine &engine = audio::AudioEngine::Instance();
    audio::AudioEngineConfig cfg;
    if (!engine.Initialize(cfg)) {
        logger.Error("[NkAudioPlayer] AudioEngine::Initialize a echoue (device indisponible ?)");
        return -2;
    }
    logger.Info("[NkAudioPlayer] device reel : {0} Hz, {1} canaux", engine.GetSampleRate(), engine.GetChannels());
\end{nkcode}

\texttt{AudioEngine} est un \textbf{singleton} : \texttt{Instance()} rend
toujours la même. \texttt{Initialize} ouvre le périphérique de sortie et lance le
fil de mixage.

La ligne de journalisation n'est pas décorative. Le périphérique impose son taux
d'échantillonnage et son nombre de canaux ; votre \texttt{cfg} n'exprime qu'un
souhait. \textbf{Lisez \texttt{GetSampleRate()} et \texttt{GetChannels()} après
\texttt{Initialize()}, jamais avant.}

\begin{nkcode}{Kernel/Runtime/NKAudio/src/NKAudio/NkAudio.h:396-411}
    struct AudioEngineConfig {
            AudioBackendType backend = AudioBackendType::AUTO;
            int32 sampleRate = AUDIO_DEFAULT_SAMPLE_RATE;   // 48000
            int32 channels   = AUDIO_DEFAULT_CHANNELS;      // 2
            int32 bufferSize = AUDIO_DEFAULT_BUFFER_SIZE;
            int32 maxVoices  = AUDIO_MAX_VOICES;
            bool enableHrtf = false;    bool enableDoppler = true;
            float32 masterVolume = 1.0f;
            memory::NkAllocator *allocator = nullptr;       ///< nullptr = allocateur global
            bool enableMasterLimiter = true;
            float32 masterLimiterThresholdDb = -0.5f;
            ResamplingQuality resamplingQuality = ResamplingQuality::LINEAR;
    };
\end{nkcode}

Les valeurs par défaut conviennent. \textbf{Ne touchez pas au champ
\texttt{backend}} — nous verrons dans un instant pourquoi.

\begin{nkretenir}[Un seul \texttt{Initialize}, un seul \texttt{Shutdown}]
\begin{nkcode}{Kernel/Runtime/NKAudio/src/NKAudio/NkAudio.h:1071-1073}
 @note Initialize() et Shutdown() sont appelés une seule fois
       depuis le thread principal. Toutes les autres méthodes
       sont thread-safe via des atomics et queues lock-free.
\end{nkcode}
Les contrats sont formels : \texttt{Initialize} exige
\texttt{!IsInitialized()}, \texttt{Shutdown} exige \texttt{IsInitialized()}
(\nkfile{NkAudio.h:1100-1111}). Toutes les \emph{autres} méthodes, elles, sont
appelables de n'importe où.
\end{nkretenir}

\subsection{Temps 2 — décoder le fichier}

\begin{nkcode}{Applications/NkAudioPlayer/src/main.cpp:115-121}
    // ── 2) Charge/decode le fichier ───────────────────────────────────────────
    audio::AudioSample sample = audio::AudioLoader::Load(path.CStr());
    if (!sample.IsValid()) {
        logger.Error("[NkAudioPlayer] impossible de charger/decoder : {0}", path.CStr());
        engine.Shutdown();
        return -3;
    }
\end{nkcode}

\texttt{AudioLoader} est entièrement statique :

\begin{nkcode}{Kernel/Runtime/NKAudio/src/NKAudio/NkAudio.h:583-648}
static AudioFormat DetectFormat(const uint8* data, usize size);
static AudioFormat DetectFormatFromPath(const char* path);
static AudioSample Load(const char* path, memory::NkAllocator* allocator = nullptr);
static AudioSample LoadFromMemory(const uint8* data, usize size,
                                  AudioFormat format = AudioFormat::UNKNOWN,
                                  memory::NkAllocator* allocator = nullptr);
static bool SaveWAV(const char* path, const AudioSample& sample);
static void Free(AudioSample& sample);
static void ConvertSampleFormat(AudioSample& sample, SampleFormat target);
static void Resample(AudioSample& sample, int32 targetSampleRate, bool highQuality = true);
static void ConvertChannels(AudioSample& sample, int32 targetChannels);
\end{nkcode}

\texttt{Load} détecte le format par le contenu, pas par l'extension : renommer un
MP3 en \texttt{.wav} ne le trompe pas.

\subsection{Temps 3 — lancer une voix}

\begin{nkcode}{Applications/NkAudioPlayer/src/main.cpp:158-162}
    // ── 4) Joue + prepare la forme d'onde ─────────────────────────────────────
    audio::VoiceParams vp;
    vp.bus = "Music";
    vp.volume = 1.0f;
    audio::AudioHandle handle = engine.Play(sample, vp);
\end{nkcode}

Une \emph{voix} est une instance de lecture. Jouer trois fois le même
\texttt{AudioSample} crée trois voix indépendantes, chacune avec son volume, sa
position de lecture et son \emph{handle}.

\begin{nkcode}{Kernel/Runtime/NKAudio/src/NKAudio/NkAudio.h:346-357}
    struct VoiceParams {
            float32 volume = 1.0f;      float32 pitch = 1.0f;    float32 pan = 0.0f;
            bool looping = false;       float32 loopStart = 0.0f; float32 loopEnd = -1.0f;
            float32 fadeInTime = 0.0f;  float32 startOffset = 0.0f;
            int32 priority = 128;
            AudioSource3D source3d;
            const char *bus = "SFX";    ///< Bus de routage (SFX/Music/Voice/UI)
    };
\end{nkcode}

Le \texttt{AudioHandle} est un simple entier enveloppé
(\nkfile{NkAudio.h:227-249}), avec un \texttt{IsValid()} et une conversion
implicite en \texttt{bool}. Il sert à piloter la voix après coup :

\begin{nkcode}{Kernel/Runtime/NKAudio/src/NKAudio/NkAudio.h:1145-1166}
void Stop(AudioHandle handle, float32 fadeOutTime = 0.0f);
void Pause(AudioHandle handle);  void Resume(AudioHandle handle);
bool IsPlaying(AudioHandle) const; bool IsPaused(AudioHandle) const; bool IsLooping(AudioHandle) const;
void SetVolume/SetPitch/SetPan/SetLooping(AudioHandle, ...);
float32 GetVolume/GetPitch/GetPan(AudioHandle) const;
float32 GetPlaybackPosition(AudioHandle) const;
void    SetPlaybackPosition(AudioHandle, float32 seconds);
\end{nkcode}

\begin{nkpiege}[\texttt{Play} peut échouer sans rien dire]
\begin{nkcode}{Kernel/Runtime/NKAudio/src/NKAudio/NkAudio.h:1125}
 @return Handle de contrôle, invalide si pool plein
\end{nkcode}
Le nombre de voix simultanées est plafonné par \texttt{maxVoices}. Passé ce
seuil, \texttt{Play} renvoie un \emph{handle} invalide et rien ne se joue. Si
vous comptez réutiliser le \emph{handle}, testez \texttt{h.IsValid()}. Si vous
n'en avez pas besoin — un son de clic, par exemple — ignorer le retour est
parfaitement légitime.
\end{nkpiege}

\subsection{Temps 4 — piloter pendant la lecture}

\begin{nkcode}{Applications/NkAudioPlayer/src/main.cpp:185-191}
                else if (k->GetKey() == NkKey::NK_SPACE) {
                    paused = !paused;
                    if (paused)
                        engine.Pause(handle);
                    else
                        engine.Resume(handle);
                }
\end{nkcode}

et la détection de fin naturelle :

\begin{nkcode}{Applications/NkAudioPlayer/src/main.cpp:254-255}
        if (!engine.IsPlaying(handle) && !paused)
            running = false; // fin naturelle
\end{nkcode}

Notez la condition composée : une voix en pause n'est pas « en train de jouer ».
Sans le \texttt{\&\& !paused}, appuyer sur la barre d'espace fermerait
l'application.

\subsection{Temps 5 — l'ordre de fermeture, qui n'est pas négociable}

\begin{nkcode}{Applications/NkAudioPlayer/src/main.cpp:258-259}
    engine.Shutdown();
    audio::AudioLoader::Free(sample);
\end{nkcode}

\begin{nkpiege}[\texttt{Shutdown()} d'abord, \texttt{Free()} ensuite. Toujours.]
La documentation de \texttt{Play} pose le contrat :

\begin{nkcode}{Kernel/Runtime/NKAudio/src/NKAudio/NkAudio.h:1123}
 @param sample  Sample source (doit rester valide pendant la lecture)
\end{nkcode}

\texttt{Play} ne copie pas vos échantillons : la voix lit \emph{directement} dans
\texttt{sample.data}, depuis le fil audio. Libérer le \texttt{AudioSample} avant
d'avoir arrêté le moteur, c'est laisser un fil temps réel lire de la mémoire
rendue au système.

Le symptôme est particulièrement déplaisant : cela ne plante pas toujours. Selon
ce qui a été réalloué à cet endroit, vous obtiendrez un grésillement, un bruit
blanc, un silence — ou un plantage, mais dix secondes plus tard, ailleurs.

\textbf{Retenez la séquence : arrêter le moteur, puis libérer les données.}
Toujours dans cet ordre, dans les cinq chemins de sortie de votre programme.
\end{nkpiege}

Et la libération elle-même a sa règle :

\begin{nkcode}{Kernel/Runtime/NKAudio/src/NKAudio/NkAudio.h:284, 620}
memory::NkAllocator *mAllocator = nullptr; ///< Allocateur responsable (jamais nullptr si IsValid())
 @note Doit être appelé avec le même allocateur utilisé pour Load()
\end{nkcode}

C'est la même règle que pour NKImage : jamais \texttt{delete}, jamais
\texttt{free}. Un \texttt{IAudioStream} qu'on n'a confié à personne se détruit
avec \texttt{memory::NkGetDefaultAllocator().Delete(stream)} — le dépôt le note
explicitement, avec la mention du plantage \texttt{c0000374} en cas de mélange
(\nkfile{Applications/NkAudioDemo/src/main.cpp:44}).

\section{Jouer un son au clic d'un bouton}

Voici l'objectif annoncé. Tout est en place ; il ne reste qu'à assembler.

\subsection{Au démarrage — une seule fois}

\begin{nkcode}{Exemple écrit pour le cours (API : NkAudio.h:1091, 1104, 595)}
#include "NKAudio/NkAudio.h"
using namespace nkentseu;

audio::AudioEngine &engine = audio::AudioEngine::Instance();
if (!engine.Initialize(audio::AudioEngineConfig{})) {
    logger.Error("[UI] audio indisponible : les sons d'interface seront muets");
    // ... mais l'application continue : le son n'est pas vital.
}
audio::AudioSample clic = audio::AudioLoader::Load("Resources/Audio/clic.wav");
if (!clic.IsValid())
    logger.Warn("[UI] son de clic introuvable");
\end{nkcode}

Remarquez la posture : l'échec de l'audio ne doit pas empêcher l'application de
démarrer. Une carte son occupée, un périphérique débranché, une machine sans
sortie : ce sont des situations normales.

\subsection{Dans la frame — quand le widget renvoie \texttt{true}}

\begin{nkcode}{Exemple écrit pour le cours (API : NkGuiWidgets.h:44, NkAudio.h:1127)}
if (nkgui::Button(ctx, "Valider")) {
    audio::VoiceParams vp;
    vp.bus = "UI";                       // bus dedie : volume UI reglable a part
    engine.Play(clic, vp);               // handle ignore : son court, « fire and forget »
    // ... l'action du bouton ...
}
\end{nkcode}

Trois lignes. C'est tout — et c'est possible parce que \texttt{Play} est
\emph{non bloquant} : il inscrit la voix dans le mixeur et rend la main
immédiatement. Votre frame ne ralentit pas.

\subsection{À la fermeture}

\begin{nkcode}{Exemple écrit pour le cours (API : NkAudio.h:1114, 622)}
engine.Shutdown();
audio::AudioLoader::Free(clic);
\end{nkcode}

\begin{nkretenir}[Pourquoi le bus « UI »]
Ranger les sons d'interface sur leur propre bus n'est pas de la coquetterie :
cela permet à l'utilisateur de les couper sans toucher à la musique, par un
simple \texttt{engine.GetBus("UI")->SetMute(true)}. Si vous les mettez sur le bus
« SFX » par défaut, vous ne pourrez plus les distinguer des bruitages du jeu.
Le coût de ce choix est nul ; le coût de sa correction plus tard ne l'est pas.
\end{nkretenir}

\section{Les bus : régler des groupes de sons ensemble}

\begin{nkcode}{Kernel/Runtime/NKAudio/src/NKAudio/NkAudio.h:1246-1251}
 Le bus Master est cree automatiquement a Initialize() avec ses
 4 sous-buses standard : SFX (default), Music, Voice, UI.

 Volume effectif d'une voix = voix.volume * bus.volume *
                              bus.parent.volume * ... * Master.volume.
\end{nkcode}

Les quatre bus existent dès \texttt{Initialize()} : vous n'avez rien à créer.
Écrire \texttt{vp.bus = "Music"} suffit.

\begin{nkcode}{Kernel/Runtime/NKAudio/src/NKAudio/NkAudio.h:1255-1276}
NkAudioBus* GetMasterBus();
NkAudioBus* GetBus(const char* name);
NkAudioBus* GetOrCreateBus(const char* name, NkAudioBus* parent = nullptr);
\end{nkcode}

\begin{nkcode}{Kernel/Runtime/NKAudio/src/NKAudio/NkAudioBus.h:53-137 (abrégé)}
static constexpr int32 MAX_EFFECTS = 8;
static constexpr int32 MAX_CHILDREN = 16;
static constexpr int32 MAX_NAME_LEN = 32;

NkAudioBus* GetParent() const noexcept;
NkAudioBus* GetChild(int32 idx) const noexcept;
NkAudioBus* FindDescendant(const char* name) noexcept;
void SetVolume(float32) noexcept;  float32 GetVolume() const noexcept;
float32 GetEffectiveVolume() const noexcept;              // produit recursif
void SetMute(bool) / IsMuted() / SetSolo(bool) / IsSoloed()
bool AddEffect(IAudioEffect*) / RemoveEffect / ClearEffects
void SetSidechainFromBus(NkAudioBus* sourceBus, float32 amount = 0.7f,
                         float32 threshold = 0.05f) noexcept;
\end{nkcode}

Un cas d'usage typique — baisser la musique sans toucher aux bruitages :

\begin{nkcode}{Exemple écrit pour le cours (API : NkAudio.h:1266, NkAudioBus.h:84)}
if (audio::NkAudioBus *bus = engine.GetBus("Music"))
    bus->SetVolume(0.35f);
\end{nkcode}

\texttt{GetBus} renvoie \texttt{nullptr} si le moteur n'est pas initialisé : le
\texttt{if} n'est pas facultatif.

Le \emph{sidechain} mérite un mot, parce qu'il résout un problème courant :
\texttt{SetSidechainFromBus(voix, 0.7f)} sur le bus Music fait automatiquement
baisser la musique quand un dialogue joue. C'est ce que font tous les jeux, et
cela tient ici en un appel.

Enfin, le fondu croisé entre deux musiques :

\begin{nkcode}{Kernel/Runtime/NKAudio/src/NKAudio/NkAudio.h:1297}
AudioHandle PlayMusicCrossfade(const AudioSample& newMusic, float32 fadeTime = 2.0f,
                               const VoiceParams& params = VoiceParams{});
\end{nkcode}

\begin{nkretenir}[Le limiteur master, et pourquoi le son ne sature jamais]
\begin{nkcode}{Kernel/Runtime/NKAudio/src/NKAudio/NkAudio.h:409-411}
            bool enableMasterLimiter = true;
            float32 masterLimiterThresholdDb = -0.5f;
\end{nkcode}
Un limiteur est actif par défaut sur la sortie. C'est pourquoi mettre un volume
supérieur à 1,0 ne produit pas de distorsion : le signal est compressé au lieu
d'être écrêté. C'est un filet de sécurité utile — mais cela veut aussi dire que
si vous poussez tous vos volumes, vous n'obtiendrez pas « plus fort », vous
obtiendrez « plus compressé ».
\end{nkretenir}

\section{Les fichiers longs : le streaming}

Charger une heure de podcast en \texttt{Float32} stéréo à 48\,kHz représente
environ 1,4~Go en mémoire. Pour tout ce qui dépasse quelques dizaines de
secondes, il faut décoder au fil de l'eau.

\begin{nkcode}{Kernel/Runtime/NKAudio/src/NKAudio/Streaming/NkAudioStream.h:42-68}
    class IAudioStream {
            virtual ~IAudioStream() = default;
            virtual int32   ReadFrames(float32* outBuf, int32 maxFrames) noexcept = 0;
            virtual bool    Seek(nk_int64 frameIdx) noexcept = 0;
            virtual nk_int64 GetFrameCount() const noexcept = 0;
            virtual int32   GetSampleRate() const noexcept = 0;
            virtual int32   GetChannels() const noexcept = 0;
            virtual bool    IsEOF() const noexcept = 0;
    };
    IAudioStream* OpenAudioStream(const char* path) noexcept;
\end{nkcode}

C'est une interface \emph{pull} : on demande des \emph{frames}, on en reçoit. Il
faut donc quelqu'un pour tirer régulièrement, et à l'avance — c'est le rôle du
lecteur de flux, qui lance un fil décodeur et remplit un tampon circulaire :

\begin{nkcode}{Kernel/Runtime/NKAudio/src/NKAudio/Streaming/NkAudioStreamPlayer.h:41-135 (abrégé)}
    class AudioStreamPlayer {
            bool Init(int32 sampleRate, int32 channels, int32 ringBufferFrames = 88200) noexcept;
            void Shutdown() noexcept;
            bool Play(IAudioStream* stream, bool loop = false) noexcept;
            void Stop() noexcept;  void Pause() noexcept;  void Resume() noexcept;
            int32 ReadFrames(float32* outBuf, int32 maxFrames) noexcept;
            bool IsPlaying() const noexcept;
            void SetVolume(float32) / GetVolume()
            void SetSpeed(float32) / GetSpeed()
            void SeekContent(float32 seconds) noexcept;
            void FlushRing() noexcept;
            bool IsActive() const noexcept;  bool IsFinished() const noexcept;
    };
\end{nkcode}

L'ouverture, telle qu'elle est écrite dans la démo :

\begin{nkcode}{Applications/NkAudioDemo/src/main.cpp:27-50 (abrégé)}
static int RunStreamingTest(const char *path) {
    IAudioStream *stream = OpenAudioStream(path);
    if (!stream) {
        logger.Error("[NkAudioDemo] OpenAudioStream echec");
        return 1;
    }
    int32 sampleRate = stream->GetSampleRate();
    int32 channels = stream->GetChannels();

    AudioStreamPlayer player;
    if (!player.Init(sampleRate, channels, 88200)) {
        memory::NkGetDefaultAllocator().Delete(stream); // voir note c0000374
        return 2;
    }
    if (!player.Play(stream, /*loop=*/false)) {
        player.Shutdown();
        return 3;
    }
\end{nkcode}

\textbf{Après \texttt{Play}, le lecteur possède le flux} : c'est lui qui le
détruira. Avant \texttt{Play} — donc dans les chemins d'erreur — c'est à vous de
le faire, avec l'allocateur du moteur.

Reste à brancher ce lecteur sur le moteur. C'est le rôle du \emph{callback}
procédural.

\section{Le \emph{callback} procédural, et le fil audio}

\begin{nkcode}{Kernel/Runtime/NKAudio/src/NKAudio/NkAudio.h:1140-1141}
using ProceduralCallback = NkFunction<void(float32* buffer, int32 frames, int32 channels)>;
AudioHandle PlayProcedural(ProceduralCallback callback, const VoiceParams& params = VoiceParams{});
\end{nkcode}

Au lieu de lire un \texttt{AudioSample} figé, la voix appelle votre fonction
chaque fois qu'elle a besoin d'échantillons. C'est ainsi qu'on branche un flux :

\begin{nkcode}{Applications/NkVideoPlayer/src/main.cpp:396-402}
                audio::VoiceParams vp;
                vp.bus = "Music";
                audioHandle = engine.PlayProcedural(
                    [&streamPlayer](float32 *buf, int32 frames, int32 /*channels*/) {
                        streamPlayer.ReadFrames(buf, frames);
                    },
                    vp);
\end{nkcode}

C'est aussi ainsi qu'on écrit un synthétiseur : rien ne vous oblige à lire un
fichier, vous pouvez calculer les échantillons.

\begin{nkpiege}[Ce \emph{callback} tourne sur le fil audio]
\begin{nkcode}{Kernel/Runtime/NKAudio/src/NKAudio/NkAudio.h:1134-1139}
 Le callback est appelé depuis le thread audio pour remplir
 les buffers à la demande. Idéal pour synthétiseur logiciel.
 @note TEMPS RÉEL : callback ne doit pas allouer/locker
\end{nkcode}

« Temps réel » veut dire ici : votre fonction dispose de quelques millisecondes,
et si elle les dépasse, l'utilisateur entend un craquement.

Sont donc interdits dans ce \emph{callback} : allouer de la mémoire, prendre un
verrou, ouvrir un fichier, journaliser, redimensionner un conteneur. Tout ce qui
peut faire attendre un fil peut faire craquer le son.

Ce que vous pouvez faire : lire des tableaux préalloués, calculer, écrire dans
le tampon fourni, manipuler des atomiques.
\end{nkpiege}

\subsection{L'ordre de destruction avec un lecteur de flux}

Ce commentaire du lecteur vidéo mérite d'être lu deux fois :

\begin{nkcode}{Applications/NkVideoPlayer/src/main.cpp:672-676}
 streamPlayer.Shutdown() joint le thread décodeur et libère le IAudioStream
 qu'il possède (ContainerAudioStream/MemoryStream/WavStream) ; l'ordre importe :
 arrêter l'engine (qui appelle le callback procédural depuis son thread audio)
 AVANT de détruire streamPlayer, pour ne jamais mixer un stream en cours de
 destruction.
\end{nkcode}

D'où la séquence complète, dans le bon ordre :

\begin{nkcode}{Exemple écrit pour le cours (invariant : NkVideoPlayer/src/main.cpp:672-676)}
engine.Shutdown();     // 1) arrete le fil audio -> plus aucun appel au callback
player.Shutdown();     // 2) joint le fil decodeur et libere le stream
\end{nkcode}

\begin{nkretenir}[La règle générale des trois fermetures]
\texttt{engine.Shutdown()} vient \textbf{toujours en premier}, parce que c'est
lui qui fait tourner le fil qui touche à tout le reste. Ensuite seulement on
détruit ce qu'il lisait : les lecteurs de flux, puis les \texttt{AudioSample}.
\end{nkretenir}

\section{Le micro}

\begin{nkcode}{Kernel/Runtime/NKAudio/src/NKAudio/NkAudioCapture.h:27-75 (abrégé)}
    struct NkCaptureDeviceInfo { /* ... */ bool isDefault = false; };
    struct NkCaptureConfig {
            int32 sampleRate = 48000;  int32 channels = 1;  int32 ringSeconds = 4;
    };
    using NkAudioInCallback = NkFunction<void(const float32* interleaved, int32 frames, int32 channels)>;

    class NkAudioCapture {
            static NkVector<NkCaptureDeviceInfo> EnumerateDevices();
            bool Open(const NkCaptureConfig& config);  void Close();
            bool Start();  void Stop();  bool IsCapturing() const;
            int32 Read(float32* out, int32 maxFrames);
            int32 Available() const;
            void SetCallback(NkAudioInCallback cb);
            int32 SampleRate() const;  int32 Channels() const;
            static bool SelfTest();
    };
\end{nkcode}

Deux modes coexistent : le \emph{pull} (\texttt{Read} depuis votre boucle) et le
\emph{push} (\texttt{SetCallback}). Le premier est plus simple et suffit
largement :

\begin{nkcode}{Exemple écrit pour le cours (API : NkAudioCapture.h:53-63)}
#include "NKAudio/NkAudioCapture.h"

audio::NkAudioCapture cap;
audio::NkCaptureConfig ccfg;      // 48000 Hz, mono, ring de 4 s
if (!cap.Open(ccfg) || !cap.Start())
    return 1;

NkVector<float32> bloc; bloc.Resize(1024);
while (enregistre) {
    const int32 n = cap.Read(bloc.Data(), 1024);   // 0 si rien de dispo
    /* ... accumuler les n frames ... */
}
cap.Stop();
cap.Close();
\end{nkcode}

\texttt{Read} qui renvoie 0 n'est pas une erreur : c'est « rien de neuf ». Le
tampon circulaire fait quatre secondes par défaut ; si vous ne lisez pas assez
souvent, les plus anciennes données sont perdues.

\texttt{NkAudioCapture} est l'une des deux seules classes du module à posséder un
auto-test (\nkfile{NkAudioCapture.cpp:930}) ; l'autre est \texttt{NkDenoiser}
(\nkfile{NkDenoiser.cpp:304}). Il n'y a d'auto-test ni pour
\texttt{AudioEngine}, ni pour \texttt{AudioLoader}, ni pour les codecs.

L'application \texttt{NkMicRecord} donne au passage le test de codecs le plus
court du dépôt :

\begin{nkcode}{Applications/NkMicRecord/src/main.cpp:76-89 (abrégé)}
    // Mode décodage : NkMicRecord --decode <in.(wav|mp3|ogg|flac|opus)> <out.wav>
    // Charge via AudioLoader (auto-détection du format, dont Ogg-Opus) et
    // réécrit en WAV 16-bit — test end-to-end des codecs NKAudio.
    if (argc >= 4 && NkString(argv[1]) == NkString("--decode")) {
        audio::AudioSample smp = audio::AudioLoader::Load(argv[2]);
        if (!smp.IsValid()) {
            printf("[DECODE] ERREUR : chargement impossible : %s\n", argv[2]);
            return 1;
        }
        const bool ok = audio::AudioLoader::SaveWAV(argv[3], smp);
\end{nkcode}

\section{Tester sans carte son}

Voici la fonctionnalité la plus utile pour apprendre — et la moins connue.

\begin{nkcode}{Kernel/Runtime/NKAudio/src/NKAudio/NkAudio.h:1321}
void RenderToBuffer(float32* outputBuffer, int32 frameCount, int32 channels = 2);
\end{nkcode}

Couplée au backend \texttt{NULL\_OUTPUT}, elle permet de faire tourner tout le
moteur — voix, bus, effets, limiteur — \emph{sans aucun périphérique}, et de
récupérer le mixage dans un tableau que vous pouvez inspecter :

\begin{nkcode}{Exemple écrit pour le cours (API : NkAudio.h:122-136, 1321)}
audio::AudioEngineConfig cfg;
cfg.backend = audio::AudioBackendType::NULL_OUTPUT;   // aucune sortie physique
engine.Initialize(cfg);
engine.Play(smp);

NkVector<float32> mix; mix.Resize(48000 * 2);         // 1 s de stereo
engine.RenderToBuffer(mix.Data(), 48000, 2);          // appel SYNCHRONE
// mix contient maintenant exactement ce qui serait sorti des haut-parleurs.
\end{nkcode}

\begin{nkretenir}
\texttt{NULL\_OUTPUT} + \texttt{RenderToBuffer} transforme l'audio en calcul pur,
donc en quelque chose de \textbf{vérifiable}. Un test qui vérifie que le mixage
n'est pas silencieux, qu'il ne dépasse pas 1,0, ou qu'un bus muet produit bien
du zéro, tourne sans matériel et sans attendre. C'est ainsi qu'il faut aborder
les exercices de ce chapitre.
\end{nkretenir}

\section{État réel du module}

\subsection{Ce qui fonctionne}

\begin{center}
\begin{tabular}{@{}ll@{}}
\toprule
\textbf{Décodage} & WAV natif, MP3, OGG Vorbis, FLAC, Opus — tous écrits dans le dépôt \\
\textbf{Sortie}   & WASAPI, CoreAudio, ALSA, AAudio, WebAudio, Null \\
\textbf{Capture}  & \texttt{NkAudioCapture} complet, avec auto-test \\
\textbf{Streaming}& \texttt{IAudioStream}, \texttt{AudioStreamPlayer}, piste audio de conteneur vidéo \\
\textbf{Mixage}   & bus hiérarchiques, sidechain, fondu croisé musique, limiteur master \\
\textbf{3D}       & atténuation, occlusion, absorption de l'air, HRTF synthétique \\
\textbf{Hors ligne} & \texttt{RenderToBuffer} \\
\bottomrule
\end{tabular}
\end{center}

Le décodage mérite une précision, parce que le dépôt se calomnie lui-même : le
commentaire de tête de \nkfile{NkAudioLoader.cpp:3} parle encore de « stubs
MP3/OGG/FLAC ». \textbf{C'est périmé.} Les quatre décodeurs sont branchés sur de
vrais codecs (\nkfile{NkAudioLoader.cpp:411-429}). Encore un cas où le
commentaire ment et le code dit vrai.

\subsection{Ce qui ne fonctionne pas}

\begin{nkpiege}[Ne jamais demander \texttt{DIRECTSOUND}]
\begin{nkcode}{Kernel/Runtime/NKAudio/src/NKAudio/NkAudioBackends.cpp:445-451}
        bool DirectSoundAudioBackend::Initialize(int32 sr, int32 ch, int32 buf) {
            mSampleRate = sr;
            mChannels = ch;
            mBufferSize = buf;
            mImpl = memory::NkGetDefaultAllocator().New<DsImpl>();
            // DirectSoundCreate8, CreateSoundBuffer, etc.
            return true;
        }
\end{nkcode}

Lisez bien la dernière ligne : \texttt{return true}. Le backend annonce le
succès sans avoir ouvert quoi que ce soit. Vous obtiendrez un moteur
« initialisé », des voix « en lecture », des positions qui avancent — et un
silence absolu, sans le moindre message d'erreur.

\textbf{Laissez \texttt{cfg.backend} sur \texttt{AUTO}}, qui choisit WASAPI sous
Windows, CoreAudio sous macOS, ALSA sous Linux. Trois autres valeurs de l'énumé\-
ration — \texttt{PULSE\_AUDIO}, \texttt{OPEN\_SL\_ES}, \texttt{CUSTOM} — n'ont
pas de classe de backend correspondante dans
\nkfile{NkAudioBackends.h} : elles existent dans l'énumération et nulle part
ailleurs.
\end{nkpiege}

\begin{nkpiege}[Le rééchantillonnage « haute qualité » est du linéaire]
\begin{nkcode}{Kernel/Runtime/NKAudio/src/NKAudio/NkAudioLoader.cpp:588-592}
        void AudioLoader::ResampleSinc(AudioSample &sample, int32 targetSampleRate) {
            // Pour l'instant, fallback linéaire
            // TODO: Implémenter Kaiser-windowed Sinc pour qualité studio
            ResampleLinear(sample, targetSampleRate);
        }
\end{nkcode}

Donc \texttt{AudioLoader::Resample(sample, taux, /*highQuality=*/true)} fait
exactement la même chose que \texttt{false}. C'est le jumeau exact du piège
\texttt{NK\_LANCZOS3} du chapitre~5 : ça ne plante pas, ça ment.

Par symétrie, les valeurs \texttt{SINC\_4} et \texttt{SINC\_8} de
\texttt{ResamplingQuality} dans \texttt{AudioEngineConfig} sont à vérifier avant
d'être présentées comme différenciantes.

La bonne nouvelle, c'est que vous n'avez normalement pas à rééchantillonner
vous-même :

\begin{nkcode}{Applications/NkAudioPlayer/src/main.cpp:15-16}
 NOTE (agent NKCode) : le mixage NKAudio convertit le TAUX du fichier vers celui du
 device (fix mixBuffer/format-device). Un lecteur n'a RIEN à faire côté rééchantillonnage.
\end{nkcode}
\end{nkpiege}

\subsection{Le reste de la surface, en survol}

Le module contient bien davantage que ce chapitre ne couvre. Pour que vous
sachiez au moins que cela existe :

\begin{itemize}
  \item \texttt{IAudioEffect} (\nkfile{NkAudio.h:438}) et les effets concrets de
        \nkfile{NkAudioEffects.h} — réverbération, écho, chorus, compresseur,
        filtres, \emph{pitch shift}\dots{} Ils s'attachent à une voix
        (\texttt{AddEffect}), à un bus, ou à la sortie
        (\texttt{AddMasterEffect}). \textbf{Leur propriété reste à l'appelant}
        (\nkfile{NkAudio.h:1230}) : l'effet doit survivre à ce à quoi il est
        attaché ;
  \item l'audio 3D : \texttt{SetSourcePosition}, \texttt{SetListenerOrientation},
        \texttt{SetSourceOcclusion}. Notez que \textbf{c'est l'application qui
        calcule l'occlusion} — « L'application fait le raycast pour calculer
        cette valeur » (\nkfile{NkAudio.h:1177-1178}). Pour le HRTF, si vous
        n'avez pas de fichier de données, \texttt{GenerateSyntheticHrtf()}
        (\nkfile{NkAudio.h:1211}) n'en demande aucun ;
  \item \texttt{AudioGenerator} (\nkfile{NkAudio.h:695}) — synthèse de formes
        d'onde, enveloppes ADSR ;
  \item \texttt{AudioAnalyzer} (\nkfile{NkAudio.h:947}) — FFT, avec un
        \texttt{FftResult} qui porte les magnitudes et une conversion
        \texttt{BinToFrequency}. De quoi dessiner un spectre en temps réel.
\end{itemize}

\section{Relier le son à l'image}

NKAudio ne dépend d'aucun module graphique : le raccordement est entièrement à
votre charge, et c'est très bien ainsi. Trois motifs reviennent :

\begin{enumerate}
  \item \textbf{Le son au clic} — le widget renvoie \texttt{true}, vous appelez
        \texttt{Play} avec \texttt{vp.bus = "UI"}. C'est ce que nous avons fait ;
  \item \textbf{La forme d'onde} — \texttt{AudioSample::data} et
        \texttt{frameCount} vous donnent tout pour pré-calculer une enveloppe
        min/max par colonne de pixels, et \texttt{GetPlaybackPosition(handle)}
        vous donne la tête de lecture. C'est exactement ce que fait
        \texttt{ComputePeaks} dans
        \nkfile{Applications/NkAudioPlayer/src/main.cpp:56-87}, avant d'envoyer
        un quad par colonne au renderer ;
  \item \textbf{Le spectre} — \texttt{AudioAnalyzer} rend un \texttt{FftResult},
        une barre par \emph{bin}.
\end{enumerate}

Le calcul de la fraction lue, pour dessiner un curseur de lecture :

\begin{nkcode}{Applications/NkAudioPlayer/src/main.cpp:212-215}
        // Position de lecture -> fraction [0..1].
        const float32 pos = engine.GetPlaybackPosition(handle);
        const float32 frac = (durSec > 0.0) ? math::NkClamp((float32)(pos / durSec), 0.0f, 1.0f) : 0.0f;
        const int32 headCol = (int32)(frac * (float32)cols);
\end{nkcode}

Le \texttt{NkClamp} n'est pas superflu : la position peut légèrement dépasser la
durée en fin de lecture, à cause de la granularité du tampon.

\begin{nkbilan}
NKAudio repose sur un invariant : \textbf{un seul format interne}. Tout ce qui
entre y est converti, ce qui supprime la combinatoire des conversions et rend le
mélange possible. Le patron est en cinq temps, et il ne se raccourcit pas. Les
\textbf{bus} permettent de régler des groupes ensemble --- effets, musique, voix
--- au lieu de retoucher chaque son. Un fichier long se \emph{diffuse} au lieu
d'être chargé. Et le rappel procédural s'exécute sur le \textbf{fil audio} : ce
qu'on y fait de lent s'entend, sous forme de coupures.
\end{nkbilan}

\section{Exercices}

\begin{nkpratique}[1 — Le son au clic, complet]
Ajoutez à votre application NKGui trois boutons et trois sons différents, tous
sur le bus « UI ». Puis ajoutez un \texttt{Checkbox} « sons d'interface » qui
appelle \texttt{engine.GetBus("UI")->SetMute(...)}.

Vérifiez ensuite trois choses : que couper le bus UI n'affecte pas un son joué
sur « SFX » ; que l'application démarre normalement si vous renommez les fichiers
audio pour les rendre introuvables ; et que la fermeture n'émet aucun bruit
parasite — c'est-à-dire que vous appelez bien \texttt{Shutdown()} avant
\texttt{Free()}.
\end{nkpratique}

\begin{nkdemo}[2 — Mesurer sans écouter]
Écrivez un programme console qui initialise le moteur en \texttt{NULL\_OUTPUT},
joue un son, et rend une seconde de mixage dans un tableau avec
\texttt{RenderToBuffer}. Calculez et affichez le maximum absolu et la moyenne
quadratique du tableau.

Recommencez en réglant le bus sur \texttt{SetVolume(0.5f)} : le maximum doit être
divisé par deux. Puis avec \texttt{SetMute(true)} : le tableau doit être
rigoureusement nul. Vous venez d'écrire un test audio automatisable, sans
matériel.
\end{nkdemo}

\begin{nkdemo}[3 — Prouver l'ordre de fermeture]
Écrivez volontairement la mauvaise séquence :
\texttt{AudioLoader::Free(sample);} \emph{puis} \texttt{engine.Shutdown();} sur
un son de plusieurs secondes en cours de lecture. Lancez plusieurs fois et notez
ce que vous entendez et ce qui se passe.

Puis remettez le bon ordre. Cet exercice n'a l'air de rien : il vous vaccine
contre une classe entière de bugs qu'on ne diagnostique pas au débogueur, parce
que le plantage arrive ailleurs et plus tard.
\end{nkdemo}

\begin{nkpratique}[4 — Un synthétiseur en dix lignes]
Utilisez \texttt{PlayProcedural} avec un \emph{callback} qui remplit le tampon
d'une sinusoïde, en gardant la phase dans une variable capturée par référence.
Ajoutez un \texttt{SliderFloat} NKGui qui règle la fréquence.

Attention : le \emph{slider} s'écrit depuis le fil principal et le
\emph{callback} lit depuis le fil audio. Réfléchissez à ce que vous partagez —
et relisez l'interdiction d'allouer et de verrouiller. Une variable atomique de
type \texttt{float} est la bonne réponse ; un \texttt{NkVector} redimensionné
dans le \emph{callback} est la mauvaise.
\end{nkpratique}

\begin{nkquestions}
\begin{enumerate}
\item Pourquoi un format interne unique plutôt que de gérer tous les formats ?
\item Citez les cinq temps du patron canonique.
\item À quoi sert un bus, et que permet-il qu'un réglage par son ne permet pas ?
\item Quand faut-il diffuser un fichier plutôt que le charger ?
\item Que ne doit-on jamais faire dans le rappel audio, et pourquoi ?
\end{enumerate}
\end{nkquestions}

\begin{nklogique}
Un son se déclenche correctement au clavier, et « ne marche pas » au clic de
souris --- alors que les deux chemins appellent la même fonction de lecture.

\begin{enumerate}
\item Le son est-il forcément en cause ? Argumentez, puis dites où vous
      chercheriez \emph{en premier}.
\item Énumérez trois causes qui n'ont rien à voir avec l'audio.
\item Proposez une mesure qui sépare « le son n'est pas joué » de « le son est
      joué et inaudible ». Précisez ce qu'elle affiche dans chaque cas.
\end{enumerate}
\end{nklogique}
